\section{Introdução}

As métricas DORA (\textit{DevOps Research and Assessment})~--- frequência de deploy, lead time, taxa de falha de mudanças (CFR, \textit{Change Failure Rate}) e tempo de recuperação~--- tornaram-se a referência para medir o desempenho de entrega de software \cite{forsgren2018accelerate}. Elas classificam equipes em quatro faixas (\textit{Elite}, \textit{High}, \textit{Medium} e \textit{Low}) e associam entrega rápida e estável a melhores resultados organizacionais \cite{dora2024report}. Apesar disso, a maior parte da evidência vem de pesquisas com questionários respondidos por profissionais, e pouco se sabe sobre como projetos open-source populares se posicionam nessas faixas quando as métricas são calculadas a partir dos próprios dados do repositório.

O problema é que, em repositórios públicos, as métricas DORA não são observadas diretamente: o GitHub não registra deploys nem falhas em produção, e é preciso inferi-los a partir de aproximações (\textit{proxies}), como releases, commits e execuções de \textit{workflows} do GitHub Actions. Cada escolha de definição (o que conta como deploy, como identificar uma falha, quando há recuperação) pode alterar a classificação de um repositório, e não há consenso sobre qual operacionalização é a mais adequada.

O objetivo deste trabalho é mensurar as quatro métricas DORA em repositórios populares do GitHub que usam GitHub Actions, em uma janela de um ano (01/10/2025 a 30/09/2026), e investigar como elas se relacionam entre si, como variam com características dos repositórios e quão sensíveis são às definições adotadas. Para isso, formulamos oito questões de pesquisa (RQ 01 a RQ 08): as RQ 01 a RQ 04 descrevem cada métrica; a RQ 05 trata da relação entre frequência de deploy e taxa de falha; a RQ 06, das características dos repositórios associadas a melhor desempenho; a RQ 07, da sensibilidade da classificação às definições; e a RQ 08, do \textit{rework rate} e da evolução temporal das métricas. As hipóteses de cada questão, escritas antes da coleta e da análise dos dados, são apresentadas a seguir.

\paragraph{RQ 01 (frequência de deploy).}
Esperamos que a maioria dos repositórios populares com GitHub Actions publique releases com frequência semanal ou menor, e que apenas uma minoria alcance a faixa \textit{Elite} (deploys sob demanda, várias vezes por semana). Muitos projetos open-source seguem ciclos de release manuais ou por marcos, e a popularidade em estrelas indica visibilidade, não cadência de entrega; a distribuição deve ser assimétrica, com poucos repositórios muito frequentes e uma cauda longa de projetos que publicam poucas vezes ao ano.

\paragraph{RQ 02 (lead time).}
Esperamos lead times medianos de dias a poucas semanas, com a maioria dos repositórios nas faixas \textit{Medium} e \textit{Low}. Como o deploy é a publicação de uma release, e releases em projetos open-source acumulam commits por semanas antes de serem cortadas, o tempo entre o commit e a release tende a ser longo. Na variante (a), esperamos valores menores, porque ela é puxada pelos commits mais recentes antes de cada release; na variante (b), esperamos valores maiores e mais dispersos, porque ela incorpora os commits mais antigos de cada intervalo. Em ambas, a distribuição deve ter cauda longa e a mediana deve ser mais informativa que a média.

\paragraph{RQ 03 (taxa de falha de mudanças, CFR).}
Esperamos CFR baixo para a maior parte dos repositórios, na faixa \textit{Elite} ou \textit{High} (até cerca de 15\%), mas com grande dispersão. Na variante (a), baseada em workflow runs com falha, esperamos valores mais altos, porque falhas de CI ocorrem antes do deploy e incluem testes instáveis e problemas de infraestrutura que nunca chegam ao usuário. Na variante (b), baseada em releases corretivas, esperamos valores menores, porque só mudanças que chegaram a ser publicadas e depois corrigidas entram na conta, e a heurística tende a perder correções feitas sem nova release.

\paragraph{RQ 04 (tempo de recuperação).}
Esperamos que a maior parte das falhas seja recuperada em menos de um dia, já que o CI é refeito a cada push e corrigir uma falha costuma exigir um commit pequeno, mas com cauda longa de episódios que duram dias ou semanas em projetos de baixa atividade ou mantidos por poucas pessoas. Por isso, esperamos mediana na faixa \textit{Elite} ou \textit{High} (menos de um dia ou até uma semana) e média bem acima da mediana, além de uma parcela de episódios censurados, sem recuperação dentro da janela.


\paragraph{RQ 05 (frequência de deploy e taxa de falha).}
Para o CFR (a), esperamos correlação de Spearman fraca ($|\rho| < 0{,}3$) com a frequência de deploy, em linha com a afirmação do DORA de que velocidade e estabilidade não são um \textit{trade-off} \cite{dora2024report}. A frequência de deploy conta releases, enquanto o CFR (a) conta execuções de CI, e a proporção de execuções com falha depende mais da qualidade dos testes, de testes instáveis e da infraestrutura do CI do que da cadência de releases. Para o CFR (b), esperamos correlação positiva e moderada, contrariando o DORA, mas em parte por construção da métrica: quem publica várias releases por semana quase sempre tem uma nova release em até 7 dias, e basta que ela seja uma correção para que a anterior conte como falha. Esperamos, portanto, que as duas variantes levem a conclusões diferentes sobre o \textit{trade-off}.

\paragraph{RQ 06 (características associadas ao desempenho).}
Esperamos as associações mais fortes com o número de contribuidores e com o tipo do projeto. Repositórios com mais contribuidores recebem mais mudanças e têm mais pessoas para corrigir falhas, o que deve aumentar a frequência de deploy e reduzir o tempo de recuperação. Bibliotecas e \textit{frameworks} publicam versões no próprio GitHub, com \textit{patches} frequentes, enquanto aplicações e serviços costumam entregar por outros canais e publicar releases mais espaçadas, o que deve se refletir na frequência de deploy e no lead time. Para popularidade, linguagem e idade, esperamos diferenças pequenas: todos os repositórios já são populares (mais de 1.000 estrelas), o que restringe a variação, e estrelas medem visibilidade, não processo de entrega. Esperamos tamanhos de efeito pequenos a médios e que poucas comparações continuem significativas após a correção de Holm.

\paragraph{RQ 07 (sensibilidade às definições).}
Esperamos que a classificação geral mude para uma parcela relevante dos repositórios (entre 20\% e 40\% em cada par de combinações), quase sempre para uma categoria vizinha, o que deve resultar em kappa de Cohen ponderado moderado a substancial. Em relação a C1, as combinações C2 e C3 devem deslocar os repositórios para categorias melhores: somar pré-releases ou usar tags aumenta a frequência de deploy, e o lead time (b) tende a ser menor que o (a), que é puxado pelo commit mais antigo de cada release. Esperamos a maior divergência no par C1$\times$C2, que muda as três definições ao mesmo tempo, e a menor no par C2$\times$C3, que compartilha a variante de lead time. Como a categoria geral é a mediana de quatro notas, parte das mudanças em uma única métrica deve ser absorvida, e a variação por métrica deve ser maior que a da classificação geral.

\paragraph{RQ 08 (\textit{rework rate} e evolução temporal).}
Para o \textit{rework rate}, esperamos mediana entre 20\% e 40\% de releases corretivas, porque, em projetos que seguem o Versionamento Semântico, versões de \textit{patch} são comuns. Esperamos que ele seja maior que o CFR (b) e fortemente correlacionado com ele: o CFR (b) atribui a falha à release que precisou de correção e só considera correções em até 7 dias, enquanto o \textit{rework rate} conta as próprias releases corretivas, inclusive as publicadas depois desse prazo e as que corrigem outra correção. Para a evolução temporal, esperamos que praticamente nenhum repositório apresente tendência significativa em qualquer direção. Um ano é pouco tempo para mudanças de processo e, com apenas quatro trimestres por repositório, o teste de Mann-Kendall bilateral não consegue atingir $p < 0{,}05$ (o menor p-valor possível é cerca de 0,09), de modo que a ausência de tendência refletirá também a falta de poder do teste.

